\cotefile{137}
\scanpage{45}

% 1–3, upper edge
produits $\prod_{i}$ aussi -- ce qui est \deleted{\unsure{vrai|même}} automatique si les $K^i$ sont nuls sauf un nb fini…]

% 4–6, upper struck paragraph
\deleted{Comme le foncteur $\operatorname{Hom}^{\bullet}_{k}(K^{\bullet},C^{\bullet})$ $(B^\circ\longrightarrow(Ab))$ dépend de façon additive des $K^i$, sa représentabilité par $K^{\bullet}$}

% 7
Définissons un hom dans $\widehat B$

% 8, (1)
\[
 F(x)\longrightarrow\operatorname{Hom}^{\bullet}_{k}(\unsure{K^{\bullet}(x)|\alpha^{\bullet}(x)},C^{\bullet}_{F})=\operatorname{Hom}^{\bullet}_{k}(K^{\bullet}(x),F(C^{\bullet})) \tag*{(1)}
\]

% 9
i.e. un hom dans $B^{\bullet}$

% 10, (1′), with brace and gloss
\[
 \underbrace{\deleted{\unsure{x}}F(x)\otimes K^{\bullet}(x)}_{\text{\unsure{représ.|repère}\ \textquotedblright}}\longrightarrow C^{\bullet}=F(C^{\bullet}) \tag*{(1$'$)}
\]

% 10, left margin, leader to underbrace
\margin{\unsure{répr.|rep.} car les $K^i(x)$ sont libres de r.f.}

% 11, (1″), with lower annotations
\[
 \underset{\substack{\rotatebox{-90}{$\in$}\\ B}}{F(x)}\otimes
 \underset{\substack{\rotatebox{-90}{$\in$}\\ Ab_k^{\bullet}}}{\operatorname{Hom}(x,C^{\bullet})}
 \underset{\mathrm{can}}{\longrightarrow}F(C^{\bullet}) \tag*{(1$''$)}
\]

% 12–13
qui provient de l'hom fonctoriel en $\deleted{\unsure{K|x}}C\in\operatorname{Ob}K$

% 14
\[
 F(x)\otimes\operatorname{Hom}(x,C)\longrightarrow F(C)
\]

% 15–20
Comme (1$''$), donc (1) est fonctoriel en $K$, et les deux membres sont des foncteurs additifs en $K$, pour prouver que \deleted{\unsure{c'est}} (1) est un iso, $\forall K$, il suffit de le voir pour $K$ de la forme \deleted{\unsure{K}} $\mathbb E^i$. Or on a alors

% 21, (2)
\[
 F(\mathbb E^i)\longrightarrow\operatorname{Hom}^{\bullet}_{k}(\mathbb E_i^{\bullet},\deleted{\unsure{\mathbb E^{\bullet}}}F(C^{\bullet})) \tag*{(2)}
\]

% 22
où on a \unsure{évidemt|évident|évidemment}

% 23
\[
 \operatorname{Hom}^{\bullet}_{k}(\mathbb E_i^{\bullet},\underset{\substack{\rotatebox{-90}{$\in$}\\Ab^{\bullet}}}{\unsure{D^{\bullet}|\mathcal C^{\bullet}}})\simeq C^i
\]

% 24–25
et on voit de suite que moyennant cet iso, (2) est l'identité, cqfd

% archive, bottom left
% archive sheet number: 44
